\documentclass[12pt]{article}
\usepackage[utf8]{inputenc}
\usepackage[T1]{fontenc}
\usepackage[margin=1in]{geometry}
\usepackage{setspace}
\usepackage{xcolor}
\usepackage[colorlinks=true, urlcolor=blue!60!black, linkcolor=blue!60!black]{hyperref}
\setlength{\parindent}{0em}
\setlength{\parskip}{0.6em}
\renewcommand{\baselinestretch}{1.0}
\begin{document}
\begin{center}
{\LARGE \textit{DHS v. Thuraissigiam}, 591 U.S. 103 (2020)\par}
\vspace{0.5em}
{\large Case Brief --- Two-Page Edition\par}
\end{center}

\section{Memory Jogger}
An asylum seeker caught 25 yards inside the border cannot use habeas to challenge a negative credible-fear finding: the Suspension Clause protects release from custody, not a new chance at admission (\href{https://supreme.justia.com/cases/federal/us/591/19-161/}{591 U.S. 103}).

\section{Detailed Case Facts}
Expedited removal under IIRIRA lets undocumented applicants avoid summary removal only by showing a credible fear of persecution, and \S\ 1252(e)(2) bars habeas review of negative findings beyond a few narrow questions. Sri Lankan national Vijayakumar Thuraissigiam was caught 25 yards inside the border in 2017; an asylum officer, a supervisor, and an immigration judge found no credible fear.

\section{Procedural History}
His habeas petition, seeking a new asylum opportunity rather than release, was dismissed; the Ninth Circuit held \S\ 1252(e)(2) unconstitutional as applied. The Supreme Court reversed.

\section{Judicial Votes}
Reversed, 7--2: Alito wrote for Roberts, Thomas, Gorsuch, and Kavanaugh; Thomas concurred; Breyer (with Ginsburg) concurred in the judgment; Sotomayor (with Kagan) dissented.

\section{Holding}
Section 1252(e)(2) violates neither the Suspension Clause nor due process as applied. Habeas in 1789 secured release from unlawful detention, not admission, and a noncitizen at the threshold of initial entry has only the admission procedures Congress provides.

\section{Analysis of Opinions}
Alito read Blackstone, Story, and the founding-era cases as ordering simple release, with any ability to remain a collateral consequence; the finality-era immigration cases rested on statute, and \textit{Boumediene} concerned release from military custody. Under \textit{Knauff} and \textit{Mezei}, a person caught just after crossing has not entered, and giving clandestine entrants more process than port-of-entry applicants would be perverse. Thomas tied ``suspension'' to detention on suspicion. Breyer would have decided only that these factual and procedural claims fall outside habeas. Sotomayor argued that habeas has always tested the legality of removal orders and that the entry fiction now lacks any limiting principle.

\section{Examples: Future Applications}
\textbf{Same-Side:} a migrant caught near the Rio Grande who seeks a new credible-fear interview gets no habeas review. \textbf{Opposite-Side:} a long-time resident or citizen detained indefinitely without bond may seek simple release. \textbf{Fence-Sitter:} expedited removal of a Chicago resident of eighteen months.

\section{Critique}
Critics call it anachronistic to demand a 1789 analogue when no immigration or asylum statutes then existed, and note that the undefined ``threshold of entry'' invites litigation and risks returning genuine refugees through unreviewed error.

\section{Key Quotations}
\begin{quote}
``For these reasons, an alien in respondent's position has only those rights regarding admission that Congress has provided by statute.''
\end{quote}
\noindent --- Justice Alito, Opinion of the Court, 591 U.S. at 140.

\end{document}
